% =====================================================================
%  Mecanum base controller -- Phase A (base_test / navigation.py)
%  Single column: the chain is narrow enough to fit.
% =====================================================================
\begin{figure}[!t]
\centering
\begin{tikzpicture}[fcchain]

% ---------------- INITIALISATION -------------------------------------
\node[fcterm] (a0) {controller start};
\node[fcproc, below=of a0] (a1)
     {get device handles\\{\tiny \texttt{wheel1..4}, \texttt{wheel1sensor..4}}};
\node[fcproc, below=of a1] (a2)
     {set position target $= \infty$\\{\tiny puts motors in velocity mode}};
\node[fcproc, below=of a2] (a3)
     {enable position sensors\\{\tiny at the simulation time step}};
\node[fcproc, below=of a3] (a4)
     {precompute $L = l_x + l_y = 0.386$};

% ---------------- MAIN LOOP ------------------------------------------
\node[fcio, below=9mm of a4] (b0) {commanded twist $(v_x, v_y, \omega_z)$};
\node[fcsub, below=of b0] (b1)
     {\texttt{body\_to\_wheel\_speeds()}\\{\tiny Eq.~(1)--(4)}};
\node[fcdec, below=8mm of b1] (b2)
     {$\max_j |\omega_j| > \omega_{\max}$?};
\node[fcproc, below=9mm of b2] (b3)
     {scale ALL four wheels\\by $\omega_{\max}/\max_j|\omega_j|$\\
      {\tiny Eq.~(5): preserves direction}};
\node[fcproc, below=of b3] (b4) {write the four wheel velocities};
\node[fcsub, below=of b4] (b5)
     {read $\Delta\theta_i$, integrate odometry\\
      {\tiny mid-point heading}};

% ---------------- TERMINATION ----------------------------------------
\node[fcproc, below=10mm of b5] (c0) {command zero velocity to all wheels};
\node[fcterm, below=of c0] (c1) {controller exit};

% ---------------- edges ----------------------------------------------
\draw[fcline] (a0) -- (a1);
\draw[fcline] (a1) -- (a2);
\draw[fcline] (a2) -- (a3);
\draw[fcline] (a3) -- (a4);
\draw[fcline] (a4) -- (b0);
\draw[fcline] (b0) -- (b1);
\draw[fcline] (b1) -- (b2);
\draw[fcline] (b2) -- node[fclab,right] {yes} (b3);
\draw[fcline] (b2.east) -- ++(11mm,0)
      node[fclab,pos=0.6,above] {no} |- (b4.east);
\draw[fcline] (b3) -- (b4);
\draw[fcline] (b4) -- (b5);
\draw[fcback] (b5.west) -- ++(-13mm,0) |- (b0.west);
\node[font=\sffamily\tiny\itshape, text=phaseMainLine, rotate=90]
     at ($(b2.west)+(-16mm,-4mm)$) {every step};
\draw[fcline] (b5) -- (c0);
\draw[fcline] (c0) -- (c1);

% ---------------- phase bands ----------------------------------------
\begin{scope}[on background layer]
  \phaseband{bandinit}{INITIALISATION}{phaseInitLine}{(a0)(a1)(a2)(a3)(a4)}{ba}
  \phaseband{bandmain}{MAIN LOOP}{phaseMainLine}{(b0)(b1)(b2)(b3)(b4)(b5)}{bb}
  \phaseband{bandend}{TERMINATION}{phaseEndLine}{(c0)(c1)}{bc}
\end{scope}

\end{tikzpicture}
\caption{Mecanum base controller (\phaseA{}). The initialisation band
holds everything computable once, including the Webots idiom of setting
a motor's position target to infinity to obtain velocity control. The
saturation test is the only branch in the loop, and it scales all four
wheels together so that a speed limit costs magnitude and never
direction. The termination band exists because a controller that exits
without zeroing the wheels leaves the base coasting into the next run.}
\label{fig:flow_base}
\end{figure}
